% Generated draft: requires mathematical and source review.
\documentclass[11pt]{article}
\usepackage[margin=25mm]{geometry}
\usepackage{fontspec}
\setmainfont{Latin Modern Roman}
\usepackage{amsmath,amssymb,mathtools,unicode-math}
\setmathfont{Latin Modern Math}
\usepackage{enumitem,needspace,xurl,hyperref}
\setlength{\parindent}{0pt}
\setlength{\parskip}{5pt}
\setlength{\emergencystretch}{3em}
\newcommand{\N}{\mathbb{N}}
\newcommand{\Z}{\mathbb{Z}}
\newcommand{\Q}{\mathbb{Q}}
\newcommand{\R}{\mathbb{R}}
\newcommand{\C}{\mathbb{C}}
\newcommand{\F}{\mathbb{F}}
\newcommand{\A}{\mathbb{A}}
\newcommand{\PP}{\mathbb{P}}
\newcommand{\E}{\mathbb{E}}
\newcommand{\eps}{\varepsilon}
\newcommand{\dd}{\,\mathrm{d}}
\begin{document}
\section*{UnsolvedMath ID: 2752}
\textbf{Source problem code:} catalogue entry
\subsection{Definitions and mathematical statement}
An orientable surface of finite type has finite genus and finitely many punctures and boundary components. In a common connected-surface convention, $\operatorname{Mod}(S)$ is the group of isotopy classes of orientation-preserving self-homeomorphisms that fix the boundary pointwise and preserve the puncture set, with isotopies satisfying the same requirements. The source does not specify connectedness, boundary conditions, whether punctures are fixed individually, or whether orientation-reversing classes are included. This definition is therefore an explicit conventional model, not a selected answer to all those scope choices.

For chosen surfaces $S_1,S_2$ and a subgroup $H\leq\operatorname{Mod}(S_1)$ of finite index, the algebraic data in the question are an injective group homomorphism
\[
\varphi:H\longrightarrow\operatorname{Mod}(S_2).
\]
Finite index means that $\operatorname{Mod}(S_1)/H$ has finitely many left cosets; it does not mean the target image has finite index. Injectivity means $\ker\varphi=\{1\}$. The parenthetical full-group case is $H=\operatorname{Mod}(S_1)$.

The research task is to identify hypotheses on the surfaces, subgroup, and mapping-class conventions under which every such $\varphi$ is induced by specified geometric operations relating the surfaces, including suitable inclusions and covering or branched-covering constructions. An inclusion can induce extension by the identity when the relevant support and boundary data make that construction well defined. A covering construction may require restriction to a liftable subgroup and compatible choices of lifts; the source does not assert that every inclusion or cover automatically gives an injection. It also does not define a complete admissible list of ``manipulations,'' their direction and composition, or the equivalence allowed between a homomorphism and a geometric realization, such as conjugation or further finite-index restriction. These must be fixed before a theorem-level classification can be stated.

The supplied literature notes report Aramayona--Souto rigidity for injections of full mapping class groups when target genus is not much larger than source genus, but give no numerical genus range or complete exceptions. They describe a range of braid-group injections announced by Chen--Kordek--Margalit and refer to Problem 2.7. They expressly distinguish both full-group settings from arbitrary finite-index subgroup maps, the harder virtual problem in the statement. No full-group theorem is silently extended to every $H$.

The source also mentions related variants replacing groups with curve graphs and other surface-associated graphs, cites Ivanov and Brendle--Margalit as central results, and records an announced model-theoretic interpretation-to-virtual-injection connection from Disarlo--Koberda--de la Nuez González. These are related programmes, not additional hypotheses or substitute conclusions for the original homomorphism classification. Their full theorem statements are not supplied.

The registry is open and the literature triage dated 2026-08-17 classifies progress as partial: genus-range rigidity is known, but no general classification of injections from arbitrary finite-index subgroups is supplied. These results, announcements, and statuses are attributed source reports, not independent literature verification in this draft.
\subsection{Short English statement}
When must an injection from a finite-index subgroup of one surface mapping class group into another come from geometric operations on the surfaces? Known full-group rigidity covers genus ranges, while the broad virtual case remains open. Surface conventions and the allowed geometric operations still need precise specification.
\subsection{Sources}
\begin{itemize}
\item K3: A New Problem List in Low-Dimensional Topology, Problem 2.4 (2026), supplied Kirby-style source: \url{https://aimath.org/pastworkshops/kirbylistrep.pdf}. This citation is retained without opening or generating a PDF.
\item [AS12] Javier Aramayona and Juan Souto, ``Homomorphisms between mapping class groups,'' Geometry and Topology 16(4) (2012), 2285--2341, DOI 10.2140/gt.2012.16.2285.
\item [CKM19] Lei Chen, Kevin Kordek, and Dan Margalit, ``Homomorphisms between braid groups'' (2019), arXiv:1910.00712, supplied identifier.
\item [Iva97] Nikolai V. Ivanov, ``Automorphism of complexes of curves and of Teichmüller spaces,'' International Mathematics Research Notices 1997(14), 651--666. The supplied DOI string has a mathematical escape in its identifier and is not silently repaired.
\item [BM19] Tara E. Brendle and Dan Margalit, ``Normal subgroups of mapping class groups and the metaconjecture of Ivanov,'' Journal of the American Mathematical Society 32(4) (2019), 1009--1070, DOI 10.1090/jams/927.
\item [DKdlNG20] Valentina Disarlo, Thomas Koberda, and J. de la Nuez González, ``The model theory of the curve graph'' (2020), arXiv:2008.10490, supplied identifier.
\item Local registry data/unsolvedmath/problems.json, permanent catalogue entry, source code catalogue entry, ``Kirby Problem 2.4''; statement, complete literature notes, and literature triage dated 2026-08-17. Dataset attribution: Ulam AI and the UnsolvedMath Contributors.
\end{itemize}
\end{document}
